\section{The EDEM Framework}
\label{sec:framework}

This section formalises the Estimated Dynamic Equilibrium Model. We
state EDEM in its general form, then identify the Dynamic Equilibrium
(DE) model of \cref{sec:experiments-de} and Walrasian equilibrium as
limiting cases. The central technical claim is that two structural
features --- max-bid clearing and per-epoch feedback of clearing
prices into agent valuations --- generate persistent positive drift
in realised prices even when individual agents' estimation errors are
zero-mean. \Cref{sec:framework-asymmetry} derives this drift in a
clean order-statistic form.

\subsection{Primitive objects}
\label{sec:framework-primitives}

A market is staged on a finite torus $G$ of homes; each home
$h \in G$ has a fixed but unobservable \emph{fair value} $v^{*}(h)$ and
a \emph{market value} $v_{t}(h) \in \mathbb{R}_{>0}$ that evolves over
discrete time $t = 0, 1, 2, \dots$. The market is populated by two
agent classes:
\begin{itemize}[leftmargin=*]
\item \textbf{Sellers} $\mathcal{S}_{t} \subset G$, each occupying one
      home and posting an ask price $a_{t}(s)$;
\item \textbf{Buyers} $\mathcal{B}_{t}$, each free to walk the torus
      and place bids on sellers it lands with.
\end{itemize}
We let $Q_{s}(t) = |\mathcal{S}_{t}|$ and $Q_{d}(t) = |\mathcal{B}_{t}|$
denote the supply and demand quantities at time $t$; both are
state variables, not parameters.

\subsection{Heterogeneous estimation}
\label{sec:framework-estimation}

Each agent is endowed with a personal \emph{estimation function}.
For an agent $i$ valuing home $h$ at time $t$:
\begin{equation}
\label{eq:estimation}
E_{i,t}(h) \;=\; v_{t}(h) \cdot \bigl(1 + \varepsilon_{i,t}\bigr),
\qquad
\varepsilon_{i,t} \sim F_{i}(\sigma_{i}),
\end{equation}
where $F_{i}$ is a zero-mean error distribution with agent-specific
dispersion $\sigma_{i}$. Concretely, $\mathbb{E}[\varepsilon_{i,t}] = 0$
and $\operatorname{Var}(\varepsilon_{i,t})$ is increasing in
$\sigma_{i}$. The DE specialisation
(\cref{sec:framework-de-special}) takes $F_{i}$ as the symmetric
uniform $\mathcal{U}[-\sigma_{i},\,+\sigma_{i}]$; in the broader
EDEM the error distribution is left unspecified beyond zero-mean and
the support condition $\Pr(\varepsilon_{i,t} > -1) = 1$ that keeps
\cref{eq:estimation} strictly positive.

\paragraph{Unit convention.} Throughout the formal model,
$\sigma_{i}$ is a dimensionless dispersion parameter expressed as a
\emph{decimal}; e.g.\ $\sigma_{i} = 0.05$ corresponds to a $\pm 5\%$
maximum estimation error. The parameter tables and figure captions
report $\sigma$ as a percentage for readability; the conversion is
implicit when these values appear in \cref{eq:estimation} and its
descendants.

The dispersion parameter $\sigma_{i}$ is the \emph{divergence of
opinion} that \citet{miller1977} treats statically. Allowing
$\sigma_{i}$ to vary across agents (estimation heterogeneity) and over
time (rising or falling market-wide divergence of opinion) is the key
generalisation that makes EDEM a dynamic theory of price formation.

\subsection{Per-tick interaction}
\label{sec:framework-interaction}

At each tick:
\begin{enumerate}[leftmargin=*,nosep]
\item Buyers move and post bids: a buyer $b$ that lands on a seller
      $s$'s home draws a bid $\beta_{b,s}^{(t)} = E_{b,t}(h(s))$ from
      \cref{eq:estimation} and presents it to $s$. Bids accumulate;
      the seller tracks the best bid received.
\item Sellers process bids: when its patience timer elapses, $s$
      offers to sign with the buyer holding the best bid. The buyer
      either commits (a sale clears) or declines (the bid is
      dropped); the commitment rule is the key behavioural primitive
      and is discussed below in \cref{sec:framework-cond2}.
\end{enumerate}

\subsection{Market-clearing functional}
\label{sec:framework-clearing}

The realised market price at time $t$ is the output of a clearing
algorithm $A$ that consumes the full market state. Let
\begin{equation}
\label{eq:state}
X_{t} \;=\;
\bigl(\mathcal{S}_{t},\,\mathcal{B}_{t},\,
       \{v_{t}(h)\}_{h \in G},\,
       \{\beta_{\bullet}^{(t)}\},\,
       \{a_{t}(s)\}_{s \in \mathcal{S}_{t}},\,
       \{\text{patience}_{t}(s)\},\,
       \mathcal{H}_{t}\bigr)
\end{equation}
denote the model's state at $t$ (agent populations, home values,
outstanding bids and asks, seller patience timers, and the history
$\mathcal{H}_{t}$ of completed sales). The realised market price is
\begin{equation}
\label{eq:price-functional}
p_{t} \;=\; A(X_{t};\,\theta),
\end{equation}
where $\theta$ collects structural parameters. The reduced state
variables $Q_{s}(t)$, $Q_{d}(t)$, and a cross-population dispersion
summary $\sigma$ are the principal --- but not the only --- quantities
that $A$ depends on; we use the abbreviated notation
$A(Q_{s}(t), Q_{d}(t), \sigma)$ when the remaining state is held
implicit. In our two specific instantiations $A$ is realised
differently:
\begin{itemize}[leftmargin=*,nosep]
\item In DE, $A$ is the rolling average of the last $W$ \emph{sale
      events} (\cref{eq:price-rolling}). Sales are discrete events
      and not every tick records one.
\item In the speculative-market EDEM variant, $A$ is implicit in the
      multiplicative update rule \cref{eq:value-update}; sales are
      not realised, and the market price is the average home value.
\end{itemize}

\subsection{Stochastic supply and demand schedules}
\label{sec:framework-schedules}

Crucially, $Q_{s}$ and $Q_{d}$ in \cref{eq:price-functional} are not
the deterministic schedules $Q_{s}(p)$ and $Q_{d}(p)$ of textbook
microeconomics. They are integer-valued random walks driven by the
price increment $\Delta p_{t-1} = p_{t-1} - p_{t-2}$ via two balancer
parameters $C_{b}, C_{s}$:
\begin{align}
\label{eq:Qs-update}
Q_{s}(t) &= Q_{s}(t-1) + Z^{s}_{t},
& \mathbb{E}\bigl[Z^{s}_{t}\bigr] &= C_{b}\,\operatorname{sgn}(\Delta p_{t-1}), \\
\label{eq:Qd-update}
Q_{d}(t) &= Q_{d}(t-1) - Z^{d}_{t},
& \mathbb{E}\bigl[Z^{d}_{t}\bigr] &= C_{s}\,\operatorname{sgn}(\Delta p_{t-1}).
\end{align}
The integer increments $Z^{s}_{t}, Z^{d}_{t}$ are required because
agent counts are integers; they realise possibly non-integer expected
counts via the deterministic-plus-Bernoulli rule of
\cref{sec:implementation-balancer}, with a population floor of one
agent per side. With $C_{b}, C_{s} > 0$ rising prices add sellers
and remove buyers (mean-reverting); with $C_{b}, C_{s} < 0$ rising
prices add buyers and remove sellers (trend-following); with
$C_{b} = C_{s} = 0$ the populations are frozen and the price evolves
under the bid-selection asymmetry alone. The simulations in
\cref{sec:experiments-edem} use $C_{s} = C_{b}$ and report the
single number $C_{b}$.

Together,
\cref{eq:estimation,eq:price-functional,eq:Qs-update,eq:Qd-update}
define a coupled stochastic process whose realisation is a sample
path. Equilibrium in the textbook sense corresponds to the steady
state $\Delta p_{t} = \Delta Q_{s}(t) = \Delta Q_{d}(t) = 0$ --- a
non-generic state that the process visits only fleetingly under most
parameter regimes, which is the central observation of the paper.

\subsection{Bid commitment (Cond.~2)}
\label{sec:framework-cond2}

A buyer $b$ offered to sign at price $\beta$ commits iff $\beta$ is
at or above the buyer's own running benchmark over its outstanding
bids:
\begin{equation}
\label{eq:cond2}
b \text{ commits} \quad\Longleftrightarrow\quad \beta \;\ge\; \frac{1}{|\mathcal{B}_{b}|}\sum_{\beta' \in \mathcal{B}_{b}} \beta',
\end{equation}
where $\mathcal{B}_{b}$ is the set of all bids $b$ has currently
outstanding (the offered $\beta$ included, so $|\mathcal{B}_{b}| \ge 1$
always). The intuition is order-statistic selection on the buyer's
side: each $\beta' \in \mathcal{B}_{b}$ is a noisy estimate of fair
value, and the bids $b$ has placed on different sellers reflect
different draws of $b$'s estimation error. Committing only when the
offered $\beta$ is at or above the buyer's running benchmark
restricts $b$'s actual purchase to the upper tail of its own bid
distribution --- the homes for which $b$'s estimate happened to
land on the optimistic side. This is the buyer-side mirror of the
seller-side max-bid selection of \cref{sec:framework-asymmetry}, and
it is the second of the two ingredients that produces the order-
statistic drift. The prior write-up of the model in
\citet{arbuzov2018de} stated this rule with the inequality reversed;
we restore the version that matches the simulation results actually
reported there (\cref{sec:implementation-cond2}).

\subsection{The order-statistic drift: an analytical aside}
\label{sec:framework-asymmetry}

The bubble mechanism that Runs~6--8 exhibit is a clean consequence
of order-statistic bias under a tractable special case. Consider a
seller $s$ that has received $n \ge 2$ bids
$\beta_{1}, \dots, \beta_{n}$ from buyers with iid zero-mean
estimation errors $\varepsilon_{i} \sim \mathcal{U}[-\sigma, +\sigma]$:
\begin{equation}
\label{eq:bid-decomp}
\beta_{i} \;=\; v_{t}(h(s)) \cdot (1 + \varepsilon_{i}).
\end{equation}
The seller-level winning bid is $\max_{i} \beta_{i}$; its expected
ratio to the home's current value is
\begin{equation}
\label{eq:ostat-bias}
\mathbb{E}\!\left[\frac{\max_{i}\beta_{i}}{v_{t}(h(s))}\right]
\;=\;
1 + \mathbb{E}\!\left[\max_{i}\varepsilon_{i}\right]
\;=\;
1 + \sigma\,\frac{n - 1}{n + 1},
\end{equation}
using the standard order-statistic identity for the maximum of $n$
iid uniform variates. The right-hand side strictly exceeds one for
all $n \ge 2$ and all $\sigma > 0$; the gap grows in both $n$
(more bidders per seller) and $\sigma$ (wider divergence of
opinion). \Cref{eq:value-update} below propagates this expected gap
multiplicatively across epochs, generating the exponential drift
visible in \cref{fig:run6-bubble}.

The clean form of \cref{eq:ostat-bias} relies on (i) a single seller
and (ii) iid bids. The empirical bubble in
\cref{sec:exp-bubble} averages a related but slightly different
quantity (the buyer-side $\min_{s}$ in \cref{eq:value-update}); we
have not derived a closed form for that case and rely on the
simulation to certify $\bar{r}_{t} > 1$ across the parameter
regimes studied. \Cref{sec:disc-bias} discusses the empirical
implication of this distinction.

\subsection{DE as a special case}
\label{sec:framework-de-special}

The Dynamic Equilibrium model of \citet{arbuzov2018de} corresponds to
the following restriction of EDEM:
\begin{itemize}[leftmargin=*,nosep]
\item Estimation: $\varepsilon \sim \mathcal{U}[-\sigma_{i}, +\sigma_{i}]$,
      with $\sigma_{i} \sim \mathcal{U}[0, \bar{\sigma}]$ across agents
      (decimal convention; \cref{sec:framework-estimation}).
\item Schedules: $Q_{s}, Q_{d}$ are restored toward linear-in-price
      targets $\hat{Q}_{s}(p) = a_{s} + b_{s} p$,
      $\hat{Q}_{d}(p) = a_{d} + b_{d} p$ once per balance period $T_{B}$.
\item Clearing: $A$ is the rolling average of the last $W$ completed
      sales. Let $k(t)$ count the completed sales up to tick $t$ and
      $P^{\mathrm{sale}}_{j}$ denote the price of the $j$-th sale; then
      \begin{equation}
      \label{eq:price-rolling}
      p_{t} \;=\; \frac{1}{m_{t}} \sum_{j=k(t)-m_{t}+1}^{k(t)} P^{\mathrm{sale}}_{j},
      \qquad m_{t} = \min(W, k(t)).
      \end{equation}
\end{itemize}
This is the formulation that \cref{sec:experiments-de} simulates. The
sale-event indexing of \cref{eq:price-rolling} matters: ticks without
completed sales contribute nothing, and the average is taken over
realised events rather than time. The implementation maintains the
window as a fixed-length deque of the last $W$ sale prices.

The speculative-market EDEM variant of \cref{sec:experiments-edem}
replaces \cref{eq:price-rolling} with a multiplicative per-epoch
update on each home's value. Define:
\begin{align*}
\mathcal{W}_{b}(t) &= \{s \in \mathcal{S}_{t} : b \text{ holds the winning bid on } s \text{ during the epoch ending at } t\}, \\
\mathcal{Y}_{t}    &= \{b : \mathcal{W}_{b}(t) \neq \varnothing\}.
\end{align*}
The per-epoch update is then
\begin{equation}
\label{eq:value-update}
v_{t+T}(h) \;=\; v_{t}(h)\,\cdot\,\bar{r}_{t} \quad \forall h \in G,
\qquad
\bar{r}_{t} \;=\;
\begin{cases}
\dfrac{1}{|\mathcal{Y}_{t}|}\,\displaystyle\sum_{b \in \mathcal{Y}_{t}} \min_{s \in \mathcal{W}_{b}(t)} \dfrac{\beta_{b,s}^{(t)}}{v_{t}(h(s))}
& \text{if } \mathcal{Y}_{t} \neq \varnothing, \\[1.6ex]
1 & \text{if } \mathcal{Y}_{t} = \varnothing.
\end{cases}
\end{equation}

The update is uniform across all homes in $G$. This is a strong
\emph{propagation assumption}: real-estate ``comp'' signals
empirically diffuse non-uniformly, weighted by neighbourhood,
amenities, and time-since-sale. Adopting a uniform multiplier keeps
the model focused on the order-statistic mechanism of
\cref{sec:framework-asymmetry} rather than on diffusion geometry; we
flag it as a candidate target for extension in \cref{sec:ext-clearing}.

\subsection{Walrasian and Miller cases as limits}
\label{sec:framework-walras-miller}

Setting $\sigma_{i} \equiv 0$ for all $i$ collapses
\cref{eq:estimation} to $E_{i,t}(h) = v_{t}(h)$: every agent agrees
on the value of every home. The order-statistic drift of
\cref{eq:ostat-bias} vanishes ($\bar{r}_{t} \equiv 1$) and the
rolling clearing in \cref{eq:price-rolling} delivers a constant
price set by $\hat{Q}_{s}(p^{*}) = \hat{Q}_{d}(p^{*})$. Under the
additional assumption that the matching process and balancer
introduce no frictions (instantaneous adjustment, costless
encounters), EDEM approaches the Walrasian fixed-point equilibrium
as a limiting case. Spatial frictions, patience, and the
finite-population integerization of \cref{eq:Qs-update,eq:Qd-update}
prevent exact equivalence at finite scales.

If we instead retain heterogeneous $\sigma_{i}$ but freeze the
dynamics ($T_{B} \to \infty$, single-period sale, no value update),
we recover a \emph{Miller-style static premium mechanism}: the
$N$ shares are assigned to the $N$ most optimistic of the
$K > N$ potential investors, so the clearing price reflects the
optimism of the marginal optimist rather than the mean valuation.
Exact equivalence to the model in \citet{miller1977} requires
short-selling restrictions and the fixed-supply assumption; the
mechanism here is the same upper-tail selection but in a different
institutional framing.

The two regimes that empirical studies of divergence of opinion have
struggled to reconcile (premium vs.\ discount;
\cref{sec:bg-debate}) appear in EDEM as two ends of a continuum
indexed by the balancer parameters $C_{b}, C_{s}$ and the
ratio of estimation error to balancer strength. \Cref{sec:experiments}
demonstrates this empirically.
